% GENERATED FILE. Edit canonical lessons, then run npm run generate:books.
% canonical-source-hash: fnv1a64:c8db70b9bdbfa8d9
% canonical-lessons: SA-C335-maya, SA-C335-vesah, SA-C335-chalam, SA-C335-mohayati, SA-C335-akarsati

\chapter{Magic and Disguise}
\label{ch:sa-magic-and-disguise}
\hlchaptermodality{335}

\begin{chapteropening}
\textbf{By the end of this chapter:} \emph{I can tell a story with magic in it.}

Complete the last lesson of chapter 335: I can tell a story with magic in it.
\end{chapteropening}

\section[\texorpdfstring{\sk{माया} (māyā)}{माया (māyā)}]{\sk{माया} (māyā) — illusion; magic}

\label{lesson:SA-C335-maya}



\begin{quote}
Before the new one: say the Sanskrit for kills, then the Sanskrit for swallows.
\end{quote}

\subsection*{You'll want to know: \sk{माया}}
\textbf{\sk{माया}} — \emph{māyā} — \textquotedblleft{}illusion; magic\textquotedblright{}.

\begin{quote}
\sk{एतत्} \sk{सर्वं} \sk{राक्षसस्य} \sk{माया} \sk{अस्ति।}
\end{quote}

\begin{quote}
\emph{etat sarvaṃ rākṣasasya māyā asti.}
\end{quote}

\begin{quote}
All this is the demon's illusion.
\end{quote}

A feminine noun.

\begin{cousinweb}
\textbf{\sk{माया}} (\emph{māyā}) MW gives \emph{māyā} first as art, wisdom, extraordinary power, and later as illusion, deception, trickery and magic, without a root in the entry, so its origin is left uncertain here (MW).
\end{cousinweb}

\begin{grammarlens}[title={What you've built}]
One more word for telling what happened.
\end{grammarlens}

\subsection*{Your turn}
\begin{itemize}
\raggedright
  \item \emph{Say it:} \emph{māyā}
  \item \emph{Say it:} \emph{māyā}, once more
  \item \emph{Say it:} \emph{etat sarvaṃ rākṣasasya māyā asti.} — All this is the demon's illusion.
  \item \emph{Recall:} say the Sanskrit for kills, then the Sanskrit for swallows, then say \emph{māyā} again
\end{itemize}

\begin{tcolorbox}[breakable,colback=teal!4,colframe=teal!35!black,title={Before you move on}]
What does \sk{माया} mean? (\textquotedblleft{}Illusion; magic\textquotedblright{}.) Say it once more.
\end{tcolorbox}
\section[\texorpdfstring{\sk{वेषः} (veṣaḥ)}{वेषः (veṣaḥ)}]{\sk{वेषः} (veṣaḥ) — a disguise; dress}

\label{lesson:SA-C335-vesah}



\begin{quote}
Before the new one: say the Sanskrit for swallows, then the Sanskrit for illusion.
\end{quote}

\subsection*{You'll want to know: \sk{वेषः}}
\textbf{\sk{वेषः}} — \emph{veṣaḥ} — \textquotedblleft{}a disguise; dress\textquotedblright{}.

\begin{quote}
\sk{राक्षसः} \sk{यात्रिकस्य} \sk{वेषं} \sk{धारयति।}
\end{quote}

\begin{quote}
\emph{rākṣasaḥ yātrikasya veṣaṃ dhārayati.}
\end{quote}

\begin{quote}
The demon takes on the disguise of a pilgrim.
\end{quote}

A masculine noun.

\begin{cousinweb}
\textbf{\sk{वेषः}} (\emph{veṣaḥ}) comes from the root √viṣ, \textquotedblleft{}to be active, work\textquotedblright{}, by MW's account; MW gives \emph{veṣa} as work, and as dress, an artificial exterior, an assumed appearance (MW).
\end{cousinweb}

\begin{grammarlens}[title={What you've built}]
One more word for telling what happened.
\end{grammarlens}

\subsection*{Your turn}
\begin{itemize}
\raggedright
  \item \emph{Say it:} \emph{veṣaḥ}
  \item \emph{Say it:} \emph{veṣaḥ}, once more
  \item \emph{Say it:} \emph{rākṣasaḥ yātrikasya veṣaṃ dhārayati.} — The demon takes on the disguise of a pilgrim.
  \item \emph{Recall:} say the Sanskrit for swallows, then the Sanskrit for illusion, then say \emph{veṣaḥ} again
\end{itemize}

\begin{tcolorbox}[breakable,colback=teal!4,colframe=teal!35!black,title={Before you move on}]
What does \sk{वेषः} mean? (\textquotedblleft{}A disguise; dress\textquotedblright{}.) Say it once more.
\end{tcolorbox}
\section[\texorpdfstring{\sk{छलम्} (chalam)}{छलम् (chalam)}]{\sk{छलम्} (chalam) — a trick; deceit}

\label{lesson:SA-C335-chalam}



\begin{quote}
Before the new one: say the Sanskrit for illusion, then the Sanskrit for a disguise.
\end{quote}

\subsection*{You'll want to know: \sk{छलम्}}
\textbf{\sk{छलम्}} — \emph{chalam} — \textquotedblleft{}a trick; deceit\textquotedblright{}.

\begin{quote}
\sk{छलेन} \sk{सः} \sk{जनान्} \sk{वञ्चयति।}
\end{quote}

\begin{quote}
\emph{chalena saḥ janān vañcayati.}
\end{quote}

\begin{quote}
By a trick he cheats people.
\end{quote}

A neuter noun.

\begin{cousinweb}
\textbf{\sk{छलम्}} (\emph{chalam}) comes from the root √skhal, \textquotedblleft{}to stumble\textquotedblright{}, as MW explains it, as in \emph{skhalati}; MW gives \emph{chala} as fraud, deceit, sham, guise, a trick (MW).
\end{cousinweb}

\begin{grammarlens}[title={What you've built}]
One more word for telling what happened.
\end{grammarlens}

\subsection*{Your turn}
\begin{itemize}
\raggedright
  \item \emph{Say it:} \emph{chalam}
  \item \emph{Say it:} \emph{chalam}, once more
  \item \emph{Say it:} \emph{chalena saḥ janān vañcayati.} — By a trick he cheats people.
  \item \emph{Recall:} say the Sanskrit for illusion, then the Sanskrit for a disguise, then say \emph{chalam} again
\end{itemize}

\begin{tcolorbox}[breakable,colback=teal!4,colframe=teal!35!black,title={Before you move on}]
What does \sk{छलम्} mean? (\textquotedblleft{}A trick; deceit\textquotedblright{}.) Say it once more.
\end{tcolorbox}
\section[\texorpdfstring{\sk{मोहयति} (mohayati)}{मोहयति (mohayati)}]{\sk{मोहयति} (mohayati) — bewilders; deludes}

\label{lesson:SA-C335-mohayati}



\begin{quote}
Before the new one: say the Sanskrit for a disguise, then the Sanskrit for a trick.
\end{quote}

\subsection*{You'll want to know: \sk{मोहयति}}
\textbf{\sk{मोहयति}} — \emph{mohayati} — \textquotedblleft{}bewilders; deludes\textquotedblright{}.

\begin{quote}
\sk{राक्षसः} \sk{मायया} \sk{जनान्} \sk{मोहयति।}
\end{quote}

\begin{quote}
\emph{rākṣasaḥ māyayā janān mohayati.}
\end{quote}

\begin{quote}
The demon bewilders people with illusion.
\end{quote}

Past (he …): \sk{सः} \sk{जनान्} \sk{मोहितवान्}, \emph{saḥ janān mohitavān}.

\begin{cousinweb}
\textbf{\sk{मोहयति}} (\emph{mohayati}) comes from the root √muh, \textquotedblleft{}to be confused, go astray\textquotedblright{}, in its causative form: to bewilder, delude (MW).
\end{cousinweb}

\begin{grammarlens}[title={What you've built}]
One more word for telling what happened.
\end{grammarlens}

\subsection*{Your turn}
\begin{itemize}
\raggedright
  \item \emph{Say it:} \emph{mohayati}
  \item \emph{Say it:} \emph{mohayati}, once more
  \item \emph{Say it:} \emph{rākṣasaḥ māyayā janān mohayati.} — The demon bewilders people with illusion.
  \item \emph{Recall:} say the Sanskrit for a disguise, then the Sanskrit for a trick, then say \emph{mohayati} again
\end{itemize}

\begin{tcolorbox}[breakable,colback=teal!4,colframe=teal!35!black,title={Before you move on}]
What does \sk{मोहयति} mean? (\textquotedblleft{}Bewilders; deludes\textquotedblright{}.) Say it once more.
\end{tcolorbox}
\section[\texorpdfstring{\sk{आकर्षति} (ākarṣati)}{आकर्षति (ākarṣati)}]{\sk{आकर्षति} (ākarṣati) — draws; attracts}

\label{lesson:SA-C335-akarsati}



\begin{quote}
Before the new one: say the Sanskrit for a trick, then the Sanskrit for bewilders.
\end{quote}

\subsection*{You'll want to know: \sk{आकर्षति}}
\textbf{\sk{आकर्षति}} — \emph{ākarṣati} — \textquotedblleft{}draws; attracts\textquotedblright{}.

\begin{quote}
\sk{मधुरं} \sk{गीतं} \sk{जनान्} \sk{आकर्षति।}
\end{quote}

\begin{quote}
\emph{madhuraṃ gītaṃ janān ākarṣati.}
\end{quote}

\begin{quote}
The sweet song attracts people.
\end{quote}

Past (it …): \sk{गीतं} \sk{जनान्} \sk{आकृष्टवत्}, \emph{gītaṃ janān ākṛṣṭavat}.

\begin{cousinweb}
\textbf{\sk{आकर्षति}} (\emph{ākarṣati}) is built from \emph{ā}, \textquotedblleft{}towards\textquotedblright{}, and the root √kṛṣ, \textquotedblleft{}to drag, pull\textquotedblright{}, as in \emph{karṣati}; MW gives \emph{ā-kṛṣ} as to draw towards oneself, attract (MW).
\end{cousinweb}

\begin{grammarlens}[title={What you've built}]
One more word for telling what happened.
\end{grammarlens}

\subsection*{Your turn}
\begin{itemize}
\raggedright
  \item \emph{Say it:} \emph{ākarṣati}
  \item \emph{Say it:} \emph{ākarṣati}, once more
  \item \emph{Say it:} \emph{madhuraṃ gītaṃ janān ākarṣati.} — The sweet song attracts people.
  \item \emph{Recall:} say the Sanskrit for a trick, then the Sanskrit for bewilders, then say \emph{ākarṣati} again
\end{itemize}

\begin{tcolorbox}[breakable,colback=teal!4,colframe=teal!35!black,title={Before you move on}]
What does \sk{आकर्षति} mean? (\textquotedblleft{}Draws; attracts\textquotedblright{}.) Say it once more.
\end{tcolorbox}
